\documentclass[12pt]{article}
\usepackage{amsmath}
\usepackage{amssymb}
\usepackage{physics}
\usepackage{graphicx}
\usepackage{hyperref}

\newcommand{\ds}{\displaystyle}

\title{Motion on the Rotating Earth}
\author{}
\date{}

\begin{document}
\maketitle
% ============================================================
\section{Rotating-Frame Preliminaries}
% ============================================================

For a reference frame rotating with constant angular velocity
$\boldsymbol{\omega}$, the acceleration measured in the rotating frame is

\begin{equation}
  \vb{a}'
  =
  \vb{A}\,\vb{a}
  -
  \boldsymbol{\omega}\times
  (\boldsymbol{\omega}\times\vb{r}')
  -
  2\boldsymbol{\omega}\times
  \frac{d\vb{r}'}{dt}.
  \label{earth:acceleration-transformation}
\end{equation}

The corresponding inertial forces are therefore the centrifugal force,

\begin{equation}
  \vb{F}_{\mathrm{cf}}
  =
  -m\boldsymbol{\omega}\times
  (\boldsymbol{\omega}\times\vb{r}'),
  \label{earth:centrifugal-force-general}
\end{equation}

and the Coriolis force,

\begin{equation}
  \vb{F}_{\mathrm{cor}}
  =
  -2m\boldsymbol{\omega}\times
  \frac{d\vb{r}'}{dt}.
  \label{earth:coriolis-force-general}
\end{equation}

These expressions, derived in the previous essay, provide the rotating-frame
dynamics used throughout the derivation of the Foucault pendulum.
% ============================================================
\section{Centrifugal and Coriolis Forces on Earth}
% ============================================================

In this section we study the centrifugal and Coriolis forces near the Earth's surface. 
To keep the analysis as transparent as possible, we model the Earth as a perfect sphere. 
Although this approximation neglects effects such as the Earth's oblateness, it captures the essential
 physics and provides a convenient setting for deriving the 
 mathematical expressions used throughout this essay.

We consider a reference frame fixed to the Earth with origin at its center. 
The $z$-axis points from the South Pole toward the North Pol 
 to the North Pole, with positive direction from South to North. Earth's angular velocity
is along the $z$ axis:

\begin{equation}
  \boldsymbol{\omega} = \omega\,\vb{k}
  \label{earth:angular-velocity}
\end{equation}

where $\omega$ is the angular velocity of the Earth equal to $7.29\times10^{-5}\,\text{s}^{-1}$.

We will be interested in the motion of a body in the vicinity of the Earth surface, and the
ideal coordinates to describe such a motion are the spherical coordinates $(\rho, \theta,
\phi)$, such that,

\begin{align}
  x &= \rho\sin\phi\cos\theta \\
  y &= \rho\sin\phi\sin\theta \\
  z &= \rho\cos\phi
\end{align}

where $\rho \geq 0$ is the radial distance from the center of the reference system,
$0 \leq \theta \leq 2\pi$ is the azimuthal angle, and $0 \leq \phi \leq \pi$ is the polar
angle. Associated with the spherical coordinates are the unit vectors $\vb{e}_r$,
$\vb{e}_\theta$, and $\vb{e}_\phi$ given by,

\begin{align}
  \vb{e}_\rho   &= \sin\phi\cos\theta\,\vb{i} + \sin\phi\sin\theta\,\vb{j} + \cos\phi\,\vb{k} \\
  \vb{e}_\theta &= -\sin\theta\,\vb{i} + \cos\theta\,\vb{j} \\
  \vb{e}_\phi   &= \cos\phi\cos\theta\,\vb{i} + \cos\phi\sin\theta\,\vb{j} - \sin\phi\,\vb{k}
\end{align}

where $\vb{i}$, $\vb{j}$, and $\vb{k}$ are the unit vectors of the rotating reference system
along the $x$, $y$, and $z$ axes, respectively.

The unit vectors define a local orthonormal basis and maintain the orientation of the
reference system:

\begin{align}
  \vb{e}_\theta \times \vb{e}_\rho &= \vb{e}_\phi \\
  \vb{e}_\phi   \times \vb{e}_\theta &= \vb{e}_\rho \\
  \vb{e}_\rho   \times \vb{e}_\phi &= \vb{e}_\theta
\end{align}

and they are related to the latitude and longitude of the Earth:

\begin{itemize}
  \item $\vb{e}_\rho$: Radial unit vector pointing outward from the center of the Earth,
        through the point under consideration.
  \item $\vb{e}_\theta$: Points tangentially along lines of latitude towards the east.
  \item $\vb{e}_\phi$: Points tangentially along lines of longitude towards the South Pole.
\end{itemize}

The position vector in spherical coordinates is,

\[
  \vb{r} = \rho\,\vb{e}_\rho
\]

and the angular velocity is represented by,

\[
  \boldsymbol{\omega} = \omega\bigl(\cos(\phi)\,\vb{e}_\rho - \sin(\phi)\,\vb{e}_\phi\bigr).
\]

and finally the velocity of the body in spherical coordinates is given by,

\[
  \dot{\vb{r}} = \dot\rho\,\vb{e}_\rho + \rho\dot\theta\sin(\phi)\,\vb{e}_\theta
    + \rho\dot\phi\,\vb{e}_\phi
\]

Since the Earth's surface is naturally described by 
spherical coordinates, expressing both the angular velocity and the particle
 velocity in the local spherical basis considerably simplifies the derivation
  of the centrifugal and Coriolis forces.

\subsection{Forces in Spherical Coordinates}

In the rotating frame the body is subject to three forces: the gravitational force, the
centrifugal force, and the Coriolis force.

\begin{itemize}
  \item The gravitational force is directed radially inward and is given by,
  \[
    \vb{F}_\text{g} = -mg\,\vb{e}_\rho
  \]

  where $m$ is the mass of the body and $g$ is the acceleration due to gravity.

  \item The centrifugal force is directed radially outward from the axis of rotation and in
        spherical coordinates it is represented by the formula:

  \begin{equation}
    \vb{F}_\text{cf} = m\omega^2\rho\bigl(\sin^2(\phi)\,\vb{e}_\rho
      + \sin(\phi)\cos(\phi)\,\vb{e}_\phi\bigr)
    \label{earth:centrifugal-force-spherical}
  \end{equation}

  \item The Coriolis force is perpendicular to the velocity of the body and is quantified by
        the formula:

  \begin{align}
    \vb{F}_\text{cor}
    &= -2m\boldsymbol{\omega}\times\dot{\vb{r}}
     = -2m\omega\bigl(\cos(\phi)\,\vb{e}_\rho - \sin(\phi)\,\vb{e}_\phi\bigr)
       \times\bigl(\dot\rho\,\vb{e}_\rho + \rho\dot\theta\sin(\phi)\,\vb{e}_\theta + \rho\dot\phi\,\vb{e}_\phi\bigr) \nonumber \\
    &= -2m\omega\bigl(\rho\dot\theta\sin(\phi)\cos(\phi)\,\vb{e}_\rho\times\vb{e}_\theta
       + (\cos(\phi)\rho\dot\phi + \sin(\phi)\dot\rho)\,\vb{e}_\rho\times\vb{e}_\phi
       - \sin^2(\phi)\rho\dot\theta\,\vb{e}_\phi\times\vb{e}_\theta\bigr) \nonumber \\
    &= -2m\omega\bigl(-\rho\dot\theta\sin(\phi)\cos(\phi)\,\vb{e}_\phi
       + (\cos(\phi)\rho\dot\phi + \sin(\phi)\dot\rho)\,\vb{e}_\theta
       - \sin^2(\phi)\rho\dot\theta\,\vb{e}_\rho\bigr)
  \end{align}

  which simplifies to,

  \begin{equation}
    \vb{F}_\text{cor} = m\omega\bigl(
      2\rho\dot\theta\sin^2(\phi)\,\vb{e}_\rho
      - 2(\dot\rho\sin(\phi) + \rho\dot\phi\cos(\phi))\,\vb{e}_\theta
      + \rho\dot\theta\sin(2\phi)\,\vb{e}_\phi\bigr)
    \label{earth:coriolis-force-spherical}
  \end{equation}
\end{itemize}

Before proceeding, it is instructive to examine the combined inertial 
force obtained by adding the centrifugal and Coriolis contributions.,

\[
  \begin{aligned}
    \vb{F} &= m\omega\rho(t)\sin^2(\phi(t))\bigl(2\dot\theta(t)+\omega\bigr)\,\vb{e}_\rho \\
    &\quad - 2m\omega\bigl(\dot\rho(t)\sin(\phi(t)) + \rho(t)\dot\phi(t)\cos(\phi(t))\bigr)\,\vb{e}_\theta \\
    &\quad + m\omega\rho(t)\sin(\phi(t))\cos(\phi(t))\bigl(2\dot\theta(t)+\omega\bigr)\,\vb{e}_\phi
  \end{aligned}
\]

One particularly interesting situation emerge from these expressions, which corresponds to a body at rest in the inertial frame. \par

A body at rest in the inertial frame appears in the rotating frame to move with angular velocity
$\dot\theta=-\omega$, while $\dot\rho=\dot\phi=0$. Under these conditions, the combined inertial force is

\[
  \vb{F} = -m\omega^2\rho\bigl(\sin^2(\phi)\,\vb{e}_\rho + \sin(\phi)\cos(\phi)\,\vb{e}_\phi\bigr)
\]

The Coriolis force is therefore equal to twice the centrifugal force
 in magnitude and points in the opposite direction. Their resultant is an 
 inward force equal to the negative of the centrifugal force. This 
resultant produces the centripetal acceleration observed in the rotating
 frame as the body appears to revolve around the Earth's axis with  velocity $-\omega$.\par

In the second case, it is readily verified that the inertial force vanishes when
$\dot\rho = 0$, $\dot\phi = 0$, and $\dot\theta = -\omega/2$. Thus, although the body
appears to rotate uniformly in the rotating frame with an angular velocity equal to half
that of the Earth (in the opposite direction), the centrifugal and Coriolis forces exactly
cancel one another.

\subsection{Lagrangian in Spherical Coordinates}

We now formulate the dynamics of the system in spherical coordinates. Under the assumption
that the gravitational acceleration $g$ is constant and directed radially inward, the
gravitational potential is

\[
V_g = mg\rho.
\]

The centrifugal force derives from the position-dependent potential

\[
V_{\mathrm{cf}}
=
-\frac{1}{2}m\omega^2\rho^2\sin^2(\phi).
\]

The Coriolis force cannot be derived from an ordinary position-dependent potential. It can,
however, be represented by the velocity-dependent generalized potential

\[
V_{\mathrm{cor}}
=
-m\omega\rho^2\dot\theta\sin^2(\phi).
\]

Indeed, for each generalized coordinate
$q_i\in\{\rho,\phi,\theta\}$, it satisfies

\[
\frac{d}{dt}
\frac{\partial V_{\mathrm{cor}}}{\partial\dot q_i}
-
\frac{\partial V_{\mathrm{cor}}}{\partial q_i}
=
\left\langle
\vb F_{\mathrm{cor}}
\,\middle|\,
\frac{\partial\vb r}{\partial q_i}
\right\rangle.
\]

Explicitly,

\begin{align}
  \frac{d}{dt}\frac{\partial V_\text{cor}}{\partial\dot\theta}
    - \frac{\partial V_\text{cor}}{\partial\theta}
    &= \left\langle\vb{F}_\text{cor}\,\middle|\,\frac{\partial\vb{r}}{\partial\theta}\right\rangle, \\
  \frac{d}{dt}\frac{\partial V_\text{cor}}{\partial\dot\phi}
    - \frac{\partial V_\text{cor}}{\partial\phi}
    &= \left\langle\vb{F}_\text{cor}\,\middle|\,\frac{\partial\vb{r}}{\partial\phi}\right\rangle, \\
  \frac{d}{dt}\frac{\partial V_\text{cor}}{\partial\dot\rho}
    - \frac{\partial V_\text{cor}}{\partial\rho}
    &= \left\langle\vb{F}_\text{cor}\,\middle|\,\frac{\partial\vb{r}}{\partial\rho}\right\rangle.
\end{align}

Similarly, the centrifugal potential satisfies

\begin{align}
  -\frac{\partial V_\text{cf}}{\partial\theta}
    &= \left\langle\vb{F}_\text{cf}\,\middle|\,\frac{\partial\vb{r}}{\partial\theta}\right\rangle, \\
  -\frac{\partial V_\text{cf}}{\partial\phi}
    &= \left\langle\vb{F}_\text{cf}\,\middle|\,\frac{\partial\vb{r}}{\partial\phi}\right\rangle, \\
  -\frac{\partial V_\text{cf}}{\partial\rho}
    &= \left\langle\vb{F}_\text{cf}\,\middle|\,\frac{\partial\vb{r}}{\partial\rho}\right\rangle.
\end{align}

The kinetic energy is

\[
T
=
\frac{1}{2}m
\left(
\dot\rho^2
+
\rho^2\dot\phi^2
+
\rho^2\dot\theta^2\sin^2(\phi)
\right).
\]

The Lagrangian is therefore

\[
L
=
T
-
V_g
-
V_{\mathrm{cf}}
-
V_{\mathrm{cor}},
\]

or, after combining the terms,

\[
L
=
\frac{1}{2}m
\left(
-2g\rho
+
\rho^2
\left(
\sin^2(\phi)(\dot\theta+\omega)^2
+
\dot\phi^2
\right)
+
\dot\rho^2
\right).
\]

The appearance of the combination $(\dot\theta+\omega)$ has a direct physical interpretation:
it is the angular velocity of the body with respect to the inertial frame.

The equations of motion follow from the Euler--Lagrange equations,

\begin{align}
  \frac{d}{dt}\frac{\partial L}{\partial\dot\rho}
    -\frac{\partial L}{\partial\rho}
    &=0, \\
  \frac{d}{dt}\frac{\partial L}{\partial\dot\phi}
    -\frac{\partial L}{\partial\phi}
    &=0, \\
  \frac{d}{dt}\frac{\partial L}{\partial\dot\theta}
    -\frac{\partial L}{\partial\theta}
    &=0.
\end{align}

Since the Lagrangian does not depend explicitly on $\theta$, this coordinate is cyclic and
its conjugate momentum is conserved,

\[
m\rho^2\sin^2(\phi)\bigl(\dot\theta+\omega\bigr)
=
\text{constant}.
\]

Expanding the Euler--Lagrange equations yields the nonlinear system

\begin{equation}
  \begin{gathered}
    g
    -\rho
      \left(
        \sin^2(\phi)(\dot\theta+\omega)^2
        +
        \dot\phi^2
      \right)
    +
    \ddot\rho
    =0,
    \\[0.4em]
    \rho
    \left(
      2\ddot\phi
      -
      \sin(2\phi)(\dot\theta+\omega)^2
    \right)
    +
    4\dot\rho\dot\phi
    =0,
    \\[0.4em]
    \rho\ddot\theta\sin(\phi)
    +
    2(\dot\theta+\omega)
    \left(
      \dot\rho\sin(\phi)
      +
      \rho\dot\phi\cos(\phi)
    \right)
    =0.
  \end{gathered}
  \label{earth:spherical-equations-of-motion}
\end{equation}

\subsection{Linearized Equations of Motion}

We consider now the motion of an object covering distances over the Earth surface much
smaller than Earth radius $R$ in any direction. We will assume the dynamics is in a
neighbourhood of the position identified by spherical coordinates $(\rho_0, \theta_0, \phi_0)$,

\[
  \vb{r}_0 = \rho_0\,\vb{e}_{\rho_0}
\]

where $\rho_0 \approx R$.

Around this position a local system of reference can be constructed by the unit vectors
$\vb{e}_{\rho_0}$, $\vb{e}_{\theta_0}$, and $\vb{e}_{\phi_0}$ as,
\begin{align}
  \vb{e}_{\rho_0}   &= \sin\phi_0\cos\theta_0\,\vb{i} + \sin\phi_0\sin\theta_0\,\vb{j} + \cos\phi_0\,\vb{k} \\
  \vb{e}_{\theta_0} &= -\sin\theta_0\,\vb{i} + \cos\theta_0\,\vb{j} \\
  \vb{e}_{\phi_0}   &= \cos\phi_0\cos\theta_0\,\vb{i} + \cos\phi_0\sin\theta_0\,\vb{j} - \sin\phi_0\,\vb{k}
\end{align}

and a point in the vicinity of the position $\vb{r}_0$ has spherical coordinates
$(\rho_0+\delta\rho,\;\theta_0+\delta\theta,\;\phi_0+\delta\phi)$ and position,

 
\[
  \vb{r}
  =
  (\rho_0+\delta\rho)\,
  \vb{e}_{\rho}
  (\theta_0+\delta\theta,\phi_0+\delta\phi).
\]


which up to first order reads,
\[
  \vb{r} = \vb{r}_0
    + \delta\rho\,\vb{e}_{\rho_0}
    + \rho_0\sin(\phi_0)\,\delta\theta\,\vb{e}_{\theta_0}
    + \rho_0\,\delta\phi\,\vb{e}_{\phi_0}.
\]

Under the assumption that $\|\vb{r} - \vb{r}_0\| \ll R$, we have for the displacements in
each direction,
\begin{align}
  \delta\rho &\ll \rho_0 \\
  \rho_0\,\delta\phi &\ll \rho_0 \quad\text{or}\quad \delta\phi \ll 1 \\
  \rho_0\sin(\phi_0)\,\delta\theta &\ll \rho_0 \quad\text{or}\quad \sin(\phi_0)\,\delta\theta \ll 1
\end{align}

We will therefore assume that the dynamics can be described by the perturbations,
\[
  \bigl(\rho_0 + \rho(t),\quad \theta_0 + \theta(t),\quad \phi_0 + \phi(t)\bigr)
\]
where at all times $\rho(t) \ll \rho_0$, $\phi(t) \ll 1$, and
$\sin(\phi_0)\,\theta(t) \ll 1$.

The equations of motion \eqref{earth:spherical-equations-of-motion} can be expressed in terms of the
displacements $(\rho(t), \theta(t), \phi(t))$ and simplified to first order as,


\begin{equation}
  \begin{alignedat}{6}
       \ddot\rho(t) &=& -g
      &+\omega^2\rho_0\sin^2(\phi_0)
      &+\omega^2\sin^2(\phi_0)\,\rho(t)
      &+\omega^2\rho_0\sin(2\phi_0)\,\phi(t)
      &+2\omega\rho_0\sin^2(\phi_0)\,\dot\theta(t) \\
    2\rho_0\ddot\phi(t) &=& \omega^2\rho_0\sin(2\phi_0)
      &+\omega^2\sin(2\phi_0)\,\rho(t)
      &+2\omega^2\rho_0\cos(2\phi_0)\,\phi(t)
      &+2\omega\rho_0\sin(2\phi_0)\,\dot\theta(t) \\
    \rho_0\sin(\phi_0)\ddot\theta(t) &=&
      &-2\omega\sin(\phi_0)\,\dot\rho(t)
      &-2\omega\rho_0\cos(\phi_0)\,\dot\phi(t)
  \end{alignedat}
  \label{earth:linearized-spherical-equations}
\end{equation}

Defining dimensional variables of length as,

\begin{align}
  x &= \rho_0\,\phi & y &= \rho_0\sin(\phi_0)\,\theta & z &= \rho
\end{align}

and defining a local coordinate system as,

\begin{align}
  \vb{e}_z &= \vb{e}_{\rho_0} &
  \vb{e}_y &= \vb{e}_{\theta_0} &
  \vb{e}_x &= \vb{e}_{\phi_0}
\end{align}

the linearized equations of motion can be expressed as,

\begin{equation}
  \begin{alignedat}{6}
    \ddot{z} &=& -g
      &+\omega^2\rho_0\sin^2(\phi_0)
      &+\omega^2\sin^2(\phi_0)\,z
      &+\omega^2\sin(2\phi_0)\,x
      &+2\omega\sin(\phi_0)\,\dot{y} \\
    2\ddot{x} &=& \omega^2\rho_0\sin(2\phi_0)
      &+\omega^2\sin(2\phi_0)\,z
      &+2\omega^2\cos(2\phi_0)\,x
      &+4\omega\cos(\phi_0)\,\dot{y} \\
    \ddot{y} &=&
      &-2\omega\sin(\phi_0)\,\dot{z}
      &-2\omega\cos(\phi_0)\,\dot{x}
  \end{alignedat}
  \label{earth:local-cartesian-equations}
\end{equation}

For the analysis of the Foucault pendulum, we retain terms only up to first order
in the Earth's angular velocity $\omega$. All centrifugal terms, which are of order
$\omega^2$, are therefore neglected, while the Coriolis terms, which are of order
$\omega$, are retained. The local equations consequently reduce to

\begin{equation}
  \begin{aligned}    
    \ddot x
      &= 2\omega\cos(\phi_0)\,\dot y, \\
    \ddot y
      &= -2\omega\sin(\phi_0)\,\dot z
         -2\omega\cos(\phi_0)\,\dot x.\\
         \ddot z
      &= -g
        +2\omega\sin(\phi_0)\,\dot y, \\
  \end{aligned}
  \label{earth:first-order-local-equations}
\end{equation}

\subsection{Foucault Pendulum}

In this section, we will derive the equations of motion for the Foucault pendulum, applying
several simplifying assumptions to streamline the process. The Foucault pendulum is a
straightforward experiment that illustrates the Earth's rotation. It consists of a heavy mass
$m$ suspended at a latitude $\phi_0$ by a long wire or rod of length $l$, free to oscillate
in any direction. Over time, the plane of the pendulum's oscillation rotates due to the
Coriolis force acting upon it. \par

To derive the equations of motion for the Foucault pendulum, we begin with the equations
\eqref{earth:first-order-local-equations} from which we can derive the total force acting on the body as,

\[
  \vb{F}_\text{total} = m
  \begin{pmatrix}  \ddot{x} \\ \ddot{y} \\ \ddot{z} \end{pmatrix}
  = m
  \begin{pmatrix}
    2\omega\cos(\phi_0)\,\dot{y} \\
    -2\omega\sin(\phi_0)\,\dot{z} - 2\omega\cos(\phi_0)\,\dot{x}\\
     -g + 2\omega\sin(\phi_0)\,\dot{y} 
  \end{pmatrix}
\]

the Coriolis component of the total force is,

\[
  \vb{F}_\text{cor} =  m
  \begin{pmatrix}
    2\omega\cos(\phi_0)\,\dot{y} \\
    -2\omega\sin(\phi_0)\,\dot{z} - 2\omega\cos(\phi_0)\,\dot{x}\\
     2\omega\sin(\phi_0)\,\dot{y} 
  \end{pmatrix} = -2 m \boldsymbol{\omega} \times {\dot{\vb{p}}}
\]

where the angular velocity $\boldsymbol{\omega}$ of the Earth is considered
in this system of reference and is given by,

\[
  \boldsymbol{\omega} = \omega(-\sin(\phi_0)\,\vb{e}_x + \cos(\phi_0)\,\vb{e}_z)
\]


and the position vector $\vb{p}$ is given by $\vb{p} = x\,\vb{e}_x + y\,\vb{e}_y + z\,\vb{e}_z$.
The Coriolis force above has the generalized potential,

\[
  V_\text{Cor} =  -m\ip{\boldsymbol{\omega}\times\vb{p}}{\dot{\vb{p}}} =
 m \omega     \left(  y\cos(\phi_0)\dot{x} -( \sin(\phi_0)z +  \cos(\phi_0)x)\dot{y} + \sin(\phi_0) y \dot{z} \right)
\]

The Lagrangian of the system, including the gravitational force reads,
\[
  \begin{aligned}
    L
    &=
    \frac{1}{2}m
    \left(
      \dot{x}^2
      +
      \dot{y}^2
      +
      \dot{z}^2
    \right)
    \\[0.4em]
    &\quad
    -m\omega
    \Bigl[
      y
      \left(
        \cos(\phi_0)\,\dot{x}
        +
        \sin(\phi_0)\,\dot{z}
      \right)
      -
      \dot{y}
      \left(
        \cos(\phi_0)\,x
        +
        \sin(\phi_0)\,z
      \right)
    \Bigr]
    -
    mgz.
  \end{aligned}
\]

Utilizing spherical coordinates centered at the origin of the reference system  , with fixed radial
coordinate $\rho=l$, we write the position of the pendulum mass as

\[
  \vb p
  =
  l\sin(\phi)\cos(\theta)\,\vb e_x
  +
  l\sin(\phi)\sin(\theta)\,\vb e_y
  -
  l\cos(\phi)\,\vb e_z.
\]

or 

\begin{align}
  x &= l\sin(\phi)\cos(\theta), \\
  y &= l\sin(\phi)\sin(\theta), \\
  z &= l\cos(\phi).
\end{align}

Here $\phi$ is the angle measured from the downward vertical, while $\theta$ is the
azimuthal angle in the local horizontal plane. \par

After some elaborate algebra, the Lagrangian in these sperical coordinates reads,

\[
  \begin{aligned}
    L ={}&
      \frac{1}{2}l^2m
      \left(
        \dot{\varphi}^{\,2}
        + \sin^2(\varphi)\dot{\theta}^{\,2}
      \right) \\
    &+ l^2m\omega
      \left[
        \cos(\phi_0)\sin^2(\varphi)
        + \frac{1}{2}\sin(\phi_0)\cos(\theta)\sin(2\varphi)
      \right]\dot{\theta} \\
    &+ l^2m\omega\sin(\phi_0)\sin(\theta)\dot{\varphi}
      + glm\cos(\varphi).
  \end{aligned}
\]
\[
  \begin{aligned}
    H ={}&
      \frac{p_\varphi^2}{2l^2m}
      + \frac{p_\theta^2}{2l^2m\sin^2(\varphi)}
      + glm\cos(\varphi) \\
    &- \omega p_\theta
      \left[
        \cos(\phi_0)
        + \sin(\phi_0)\cos(\theta)
          \frac{\cos(\varphi)}{\sin(\varphi)}
      \right]
      - \omega p_\varphi\sin(\phi_0)\sin(\theta) \\
    &+ \frac{1}{2}l^2m\omega^2
      \Bigl[
        \sin^2(\phi_0)
        \left(
          \cos^2(\theta)\cos^2(\varphi)
          + \sin^2(\theta)
        \right)
        + \cos^2(\phi_0)\sin^2(\varphi) \\
    &\hspace{8em}
        + \frac{1}{2}\cos(\theta)
          \sin(2\phi_0)\sin(2\varphi)
      \Bigr].
  \end{aligned}
\].
The Euler equations of motion derived from the Lagrangian above read,

\[
  \begin{aligned}
    g\sin(\varphi)
    + l\dot{\theta}\Bigl(
        -2\omega\sin(\phi_0)\cos(\theta)\sin^2(\varphi)
        + \omega\cos(\phi_0)\sin(2\varphi)
        + \sin(\varphi)\cos(\varphi)\dot{\theta}
      \Bigr)
    - l\ddot{\varphi}
    &= 0, \\
    \sin(\varphi)\Bigl\{
      \sin(\varphi)\ddot{\theta}
      + 2\Bigl[
          \cos(\varphi)\bigl(\omega\cos(\phi_0)+\dot{\theta}\bigr)
          - \omega\sin(\phi_0)\cos(\theta)\sin(\varphi)
        \Bigr]\dot{\varphi}
      \Bigr\}
    &= 0.
  \end{aligned}
\]

which are equivalent to,

\begin{align}
  \ddot\phi(t) &= \left(-\frac{g}{l} - \omega^2\cos^2(\phi_0)\right)\phi(t) \\
  \dot\theta(t) &= -\omega\cos(\phi_0).
\end{align}

The equation in $\phi(t)$ is a simple harmonic oscillator with a frequency of
$\omega_\phi = \sqrt{g/l}$, if the higher order terms in $\omega^2\cos^2(\phi_0)$ are
removed. The equation of motion for $\theta(t)$ is a constant angular velocity which
represents the precession of the Foucault pendulum. In the northern hemisphere, the
precession is counterclockwise as $\phi < \pi/2$, while in the southern hemisphere it is
clockwise. The period of the precession is given by,
\[
  \tau = \frac{2\pi}{\omega|\cos(\phi_0)|}
    = \frac{24}{|\cos(\phi_0)|}\,\text{h}
    = \frac{24}{\sin(\phi'_0)}\,\text{h}
\]
where the last equation is obtained by substituting $\omega = 2\pi/24\,\text{h}$, and
$\phi'_0$ is the geographical latitude of the pendulum. The maximum value of the period is
at the North and South poles; it decreases when the pendulum approaches the equator where the
period is infinite.

\end{document}
